% Part: incompleteness
% Chapter: arithmetization-syntax
% Section: proofs-in-ax

\documentclass[../../../include/open-logic-section]{subfiles}

\begin{document}

\olfileid[fr]{inc}{art}{pax}
\olsection{\usetoken{P}{derivation} axiomatiques}

\begin{explain}
Pour arithmétiser les !!{derivation}s axiomatiques, nous devons les
représenter par des nombres. Comme ces !!{derivation}s sont de
simples suites de !!{formula}s, il est naturel de coder chacune
d'elles par le code de la suite des codes des !!{formula}s qu'elle
contient. Les tests de cette section portent sur les dérivations
dont toutes les lignes sont des énoncés.
\end{explain}

\begin{defn}
Si $\delta$ est une !!{derivation} axiomatique formée des
!!{formula}s $!A_1$, \dots,~$!A_n$, alors son numéro de G\"odel
$\Gn{\delta}$ est
\[
\tuple{\Gn{!A_1}, \dots, \Gn{!A_n}}.
\]
\end{defn}

\begin{ex}
  Considérons la !!{derivation} très simple :
  \begin{derivation}
  1. & $!B \lif (!B \lor !A)$ \\
  2. & $(!B \lif (!B \lor !A)) \lif (!A  \lif (!B \lif (!B \lor !A)))$\\
  3. & $!A  \lif (!B \lif (!B \lor !A))$
  \end{derivation}
  Le numéro de G\"odel de cette !!{derivation} est
  \begin{align*}
  \openTuple\,
    & \Gn{!B \lif (!B \lor !A)}, \\
    &\Gn{(!B \lif (!B \lor !A)) \lif (!A  \lif (!B \lif (!B \lor !A)))},\\
    & \Gn{!A  \lif (!B \lif (!B \lor !A))} \,\closeTuple.
  \end{align*}
\end{ex}

\begin{explain}
Une fois cette représentation des !!{derivation}s fixée, nous
devons montrer que nous pouvons les manipuler par des opérations
récursives primitives et exprimer ainsi leurs propriétés et
relations essentielles. Certaines opérations sont simples : si $d$
est le numéro de G\"odel d'!!a{derivation}, alors
$(d)_{\len{d}-1}$ est le numéro de G\"odel de sa !!{formula}
finale. D'autres propriétés sont beaucoup plus difficiles à
traiter; nous en esquisserons au moins la méthode. Notre objectif
est de montrer que la relation « $\delta$ est !!a{derivation}
de~$!A$ à partir de~$\Gamma$ » est récursive primitive sur les
numéros de G\"odel de $\delta$ et de~$!A$.
\end{explain}

\begin{prop}
\ollabel{prop:followsby}
Les relations suivantes sont récursives primitives :
\begin{enumerate}
\item $!A$ est un axiome.
\item La ligne~$i$ de~$\delta$ est justifiée par modus ponens.
\item La ligne~$i$ de~$\delta$ est justifiée par \QR.
\item $\delta$ est une !!{derivation} correcte.
\end{enumerate}
\end{prop}

\begin{proof}
Nous devons montrer que les relations correspondantes entre les
numéros de G\"odel des !!{formula}s et des !!{derivation}s sont
récursives primitives.
\begin{enumerate}
\item Nous disposons d'une liste de schémas d'axiomes; $!A$ est un
  axiome si elle est de la forme prescrite par l'un de ces schémas.
  La liste étant finie, il suffit de montrer que, pour chaque
  schéma, nous pouvons déterminer par un test récursif primitif si
  $!A$ en est une instance. Considérons par exemple le schéma
  \[
  !B \lif (!C \lif !B).
  \]
  La !!{formula}~$!A$ en est une instance s'il existe des
  !!{formula}s $!B$ et $!C$ telles que nous obtenions $!A$ en
  concaténant `$($', $!B$, `$\lif$', `$($', $!C$, `$\lif$',
  $!B$ et~`$))$'. Nous pouvons tester cette propriété sur le
  numéro de G\"odel~$n$ de~$!A$ : la concaténation des suites
  est récursive primitive et les numéros de G\"odel de $!B$ et
  de $!C$ sont inférieurs à celui de~$!A$ lorsqu'il s'agit de
  sous-!!{formula}s de~$!A$. Nous définissons donc
  \begin{multline*}
  \fn{IsAx}_{!B \lif (!C \lif !B)}(n) \defiff \bexists{b< n}{
    \bexists{c < n}{(\fn{Sent}(b) \land \fn{Sent}(c) \land {}}}\\ n =
  \Gn{(} \concat b \concat \Gn{\lif} \concat \Gn{(} \concat c \concat
  \Gn{\lif} \concat b \concat \Gn{))}).
  \end{multline*}
  Une définition analogue pour chaque schéma d'axiomes permet,
  par disjonction, de définir $\fn{IsAx}(n)$ : « $n$ est le
  numéro de G\"odel d'un axiome ».
\item La ligne~$i$ de~$\delta$ est justifiée par modus ponens si
  et seulement s'il existe deux lignes antérieures $j,k<i$ telles
  que la ligne~$j$ porte une !!{formula}~$!A$, la ligne~$k$
  porte $!A \lif !B$ et la ligne~$i$ porte~$!B$.
  \begin{multline*}
    \fn{MP}(d, i) \defiff \bexists{j < i}{\bexists{k < i}{}}\\
    (d)_k = \Gn{(} \concat (d)_j
      \concat \Gn{\lif} \concat (d)_i \concat \Gn{)}
  \end{multline*}
  La quantification bornée, la concaténation et l'égalité étant
  récursives primitives, cette formule définit une relation
  récursive primitive.
\item Une ligne de~$\delta$ est justifiée par \QR{} si elle est de
  la forme $!B \lif \lforall[x][!A(x)]$, si une ligne antérieure
  porte $!B \lif !A(c)$ pour une !!{constant}~$c$ et si $c$
  n'apparaît ni dans~$!B$ ni dans~$!A(x)$. Cela revient aux conditions suivantes :
  \begin{enumerate}
  \item il existe !!a{sentence}~$!B$ et
  \item !!a{formula}~$!A(x)$ dont $x$ est la seule variable
    éventuellement libre, tels que
  \item la ligne~$i$ porte $!B \lif \lforall[x][!A(x)]$,
  \item une ligne~$j<i$ porte $!B \lif \Subst{!A}{c}{x}$ pour une
    constante~$c$,
  \item et $c$ n'apparaît ni dans~$!B$ ni dans~$!A(x)$.
  \end{enumerate}
  Toutes ces conditions sont récursives primitives : les numéros
  de G\"odel de $!B$, de $!A(x)$ et de $x$ sont inférieurs à celui
  de la !!{formula} de la ligne~$i$. Pour borner le code de $c$,
  posons $N(d)=\sum_{r<\len{d}}\len{(d)_r}$ et
  $K(d,j)=1+\max((d)_j,\Gn{\Obj c_{N(d)}})$. Si $x$ a une occurrence
  libre dans $!A(x)$, la constante choisie figure à la ligne~$j$ et son
  numéro de G\"odel est inférieur à $(d)_j$. Sinon, la substitution
  laisse $!A(x)$ inchangée. Parmi les $N(d)+1$ premières constantes,
  l'une au moins est absente de toutes les lignes, puisque celles-ci
  contiennent au total $N(d)$ symboles. Dans les deux cas, une
  constante convenable, si elle existe, a un code inférieur à $K(d,j)$.
  On obtient
  \begin{multline*}
    \fn{QR}_1(d, i) \defiff \bexists{j < i}{\bexists{b < (d)_i}{\bexists{x <
        (d)_i}{\bexists{a < (d)_i}{\bexists{c < K(d,j)}{(}}}}}
    \\ \fn{Var}(x) \land \fn{Const}(c) \land {} \\
    (d)_i = \Gn{(} \concat
    b \concat \Gn{\lif} \concat \Gn{\lforall} \concat x \concat a
    \concat \Gn{)} \land {}\\
    (d)_j = \Gn{(} \concat b \concat
    \Gn{\lif} \concat \fn{Subst}(a,c,x) \concat \Gn{)} \land {}\\ \fn{Sent}(b)
    \land \fn{Sent}(\fn{Subst}(a,c,x)) \land {} \bforall{k <
      \len{b}}{(b)_k \neq (c)_0}\land{}\\ \fn{Frm}(a)\land\bforall{k<\len{a}}{(a)_k\neq(c)_0})
  \end{multline*}
  Nous supposons ici que $c$ et $x$ sont les numéros de G\"odel
  de la constante et de la variable considérées comme des termes,
  et non les codes de leurs seuls symboles. Pour vérifier que $x$
  est la seule variable qui puisse être libre dans~$!A(x)$, nous
  testons si $\Subst{!A(x)}{c}{x}$ est !!a{sentence}; nous imposons
  aussi que le code de chaque symbole de~$!B$ diffère du code
  $(c)_0$ du symbole de~$c$.

  Nous laissons en exercice l'autre version de~\QR.
\item $d$ est le numéro de G\"odel d'!!a{derivation} correcte si et
  seulement si chacune de ses lignes est un axiome ou se justifie
  par modus ponens ou par~\QR. Ainsi,
  \begin{multline*}
  \fn{Deriv}(d) \defiff \fn{SeqCode}(d)\land0<\len{d}\land{}\\
  \bforall{i<\len{d}}{\bigl(\fn{Sent}((d)_i)\land
    (\fn{IsAx}((d)_i)\lor\fn{MP}(d,i)\lor\fn{QR}(d,i))\bigr)}
  \end{multline*}
\end{enumerate}
\end{proof}

\begin{prob}
Définissez les relations suivantes sur le modèle de la
\olref[inc][art][pax]{prop:followsby} :
\begin{enumerate}
\item $\fn{IsAx}_{!A \lif (!B \lif (!A \land !B))}(n)$,
\item $\fn{IsAx}_{\lforall[x][!A(x)] \lif !A(t)}(n)$,
\item $\fn{QR}_{2}(d, i)$, pour l'autre version de~\QR.
\end{enumerate}
\end{prob}

\begin{prop}
Supposons que $\Gamma$ soit un ensemble récursif primitif
d'!!{sentence}s. Alors la relation $\Prf[\Gamma](x, y)$, qui
exprime « $x$ est le code d'!!a{derivation}~$\delta$ de~$!A$ à
partir de~$\Gamma$ et $y$ est le numéro de G\"odel de~$!A$ »,
est récursive primitive.
\end{prop}

\begin{proof}
Supposons que l'appartenance $y \in \Gamma$ soit exprimée par le
prédicat récursif primitif~$R_\Gamma(y)$. Nous devons montrer que
$\Prf[\Gamma](x, y)$ est récursive primitive, avec $y$ numéro de
G\"odel d'!!a{sentence}~$!A$ et $x$ code d'!!a{derivation}
de~$!A$ à partir de~$\Gamma$.

Exigeons que $y$ code un énoncé et que $x$ code une suite
non vide d'énoncés dont le dernier a pour code~$y$. On cherche
un ensemble $I$ de positions dans cette suite. Les ensembles
possibles se codent par les masques binaires inférieurs à
$2^{\len{x}}$. Pour $i\in I$, on vérifie $R_\Gamma((x)_i)$.
Pour $i\notin I$, on vérifie que la ligne est un axiome ou
qu'elle est justifiée par modus ponens ou par~\QR{} à partir de
lignes antérieures. Dans les deux cas de~\QR, on impose à
la constante témoin d'être absente de toutes les lignes
dont la position appartient à~$I$. Le même $I$ sert à toutes
les vérifications. Les recherches bornées et les tests
récursifs primitifs de la proposition précédente donnent ainsi
un test récursif primitif de $\Prf[\Gamma](x,y)$.
\end{proof}

\begin{explain}
Le théorème de la déduction donne une autre manière de coder
les témoins de la dérivabilité à partir de~$\Gamma$. On a
$\Gamma \Proves !A$ si et seulement s'il existe une liste finie
d'!!{sentence}s $!B_1$, \dots, $!B_n$ appartenant à~$\Gamma$
telle que $\{!B_1, \dots, !B_n\} \Proves !A$. Par le théorème de
déduction, cela équivaut à
$\Proves (!B_1 \lif (!B_2 \lif \cdots (!B_n \lif !A)\cdots))$.
On peut vérifier par un test récursif primitif si un énoncé de
numéro de G\"odel~$z$ possède cette forme. Dans
$\fn{Cert}_{\Gamma}(x,y)$, $x$ est le code d'une !!{derivation}
du conditionnel emboîté à partir de~$\emptyset$ et $y$ est
le code de son conséquent final.

À partir d'une suite d'!!{sentence}s et d'un !!{sentence} destiné
à être le conséquent, la récurrence primitive suivante construit
un conditionnel qui prend tous les éléments de la suite pour
antécédents :
\begin{align*}
  \fn{hCond}(s, y, 0) & = y \\
  \fn{hCond}(s, y, n+1) & = \Gn{(} \concat (s)_{n} \concat \Gn{\lif}
  \concat \fn{hCond}(s, y, n) \concat \Gn{)}\\
  \fn{Cond}(s, y) & = \fn{hCond}(s, y, \len{s})\\
  \intertext{La récurrence place les antécédents dans l'ordre inverse de la liste. On pose}
  \fn{Cert}_{\Gamma}(x,y) & \defiff \fn{Sent}(y)\land\fn{Deriv}(x)\land{} \\
  &\quad\bexists{s < \fn{sequenceBound}(x,x)}{(} \\
  &\qquad (x)_{\len{x}-1} = \fn{Cond}(s,y) \land {} \\
  &\qquad\fn{SeqCode}(s)\land\bforall{i<\len{s}}{R_\Gamma((s)_i)}).
\end{align*}
La borne portant sur~$s$ vient du fait que chaque $(s)_i$ est le
numéro de G\"odel d'une sous-!!{formula} de la dernière ligne de
la !!{derivation} et est donc inférieur à $(x)_{\len{x}-1}$.
Le nombre d'antécédents $!B$ appartenant à~$\Gamma$, autrement
dit la longueur de~$s$, est inférieur à la longueur de la
!!{formula} inscrite sur la dernière ligne de la !!{derivation}
codée par~$x$. On peut donc prendre $\fn{sequenceBound}(x,x)$
comme borne stricte : les codes des antécédents sont inférieurs
à~$x$ et leur nombre est inférieur à~$x$.
Si $y=\Gn{!A}$ pour un énoncé~$!A$, l'existence d'un
$x$ satisfaisant $\fn{Cert}_{\Gamma}(x,y)$ équivaut à
$\Gamma \Proves !A$. Dans un sens, on applique le modus ponens
successivement avec les antécédents appartenant à~$\Gamma$.
Dans l'autre, on applique le théorème de la déduction aux
hypothèses de la dérivation, en choisissant l'ordre de la liste~$s$
pour produire le conditionnel obtenu.
\end{explain}

\end{document}
